\chapter{Introduction about the IT Department}

The Information Technology (IT) Department at National Aluminium Company Limited (NALCO) plays a crucial role in managing and supporting the organization’s digital infrastructure, communication systems, enterprise applications, and industrial information systems. As a large integrated public sector enterprise with operations spread across mines, refinery units, smelter plants, power plants, and corporate offices, NALCO relies heavily on robust IT systems to ensure smooth coordination, secure data management, and efficient industrial operations.

The IT Department is responsible for maintaining the organization’s computer networks, servers, databases, enterprise software systems, and cybersecurity infrastructure. It supports the automation and digitalization of various administrative, operational, and industrial processes across different units of the organization. The department ensures uninterrupted connectivity and communication between operational locations such as Damanjodi, Angul, Bhubaneswar, and other regional offices.

One of the major responsibilities of the IT Department is the management of enterprise-level software systems used for finance, human resources, inventory management, procurement, payroll, production monitoring, and reporting operations. These systems help streamline organizational workflows, improve operational efficiency, and support data-driven decision making across departments.

The department also manages network infrastructure including Local Area Networks (LAN), Wide Area Networks (WAN), internet services, firewall systems, VPN connectivity, and secure communication channels. Regular monitoring, maintenance, troubleshooting, and system updates are carried out to ensure network reliability and data security within the organization.

Another important function of the IT Department is data management and backup administration. Large volumes of operational and administrative data generated across different NALCO units are securely stored, monitored, and maintained through centralized database systems and backup procedures. The department follows standard data security protocols, access control mechanisms, and disaster recovery procedures to protect critical organizational information from unauthorized access or system failures.

The IT Department also plays a significant role in maintaining industrial automation support systems and process monitoring applications used in refinery and mining operations. These systems assist in monitoring plant operations, production activities, equipment status, and process parameters in real time, thereby improving operational control and efficiency.

\section{Data Centre Facilities}

NALCO’s IT infrastructure includes dedicated data centre facilities that support centralized storage, processing, and management of organizational data and applications. The data centre contains servers, storage devices, networking equipment, backup systems, and security mechanisms required for running enterprise applications and maintaining uninterrupted digital operations.

The data centre environment is designed with proper cooling systems, uninterrupted power supply (UPS), backup generators, and physical security controls to ensure continuous system availability and protection of critical hardware infrastructure. The servers hosted within the data centre support database management systems, enterprise resource planning (ERP) applications, internal portals, email services, and other operational software used across the organization.

Regular data backup procedures, server monitoring, software updates, and cybersecurity measures are implemented within the data centre to ensure secure and reliable system performance. Disaster recovery mechanisms and redundancy systems are also maintained to minimize downtime and ensure business continuity during technical failures or emergencies.

\section{Procedures and Workflow Management}

The IT Department follows structured procedures for managing technical support, software maintenance, user access control, and system administration. User requests and technical issues are handled through organized support and troubleshooting processes to ensure minimal disruption in organizational operations.

Software updates, security patches, system maintenance, and database monitoring are performed regularly to maintain system stability and security. Access permissions are managed based on user roles and departmental requirements to ensure controlled access to sensitive organizational data and applications.

The department also coordinates with different operational units and departments to implement new technologies, improve digital infrastructure, and support modernization initiatives within the organization. Through continuous monitoring, maintenance, and technological support, the IT Department contributes significantly to the efficient functioning of NALCO’s industrial and administrative operations.

\section{In-House Mobile Applications and Portals}

The IT and System Department at NALCO has developed a robust suite of in-house mobile applications and digital portals to streamline enterprise workflows, enhance employee engagement, and improve customer relations. These applications ensure that critical operational data and services are securely accessible on the go.

\begin{itemize}
    \item \textbf{NALCO Netra:} A unified mobile application designed to provide active employees with convenient access to key information and self-service transactions, including payroll, attendance, and leave management.
    
    \item \textbf{NALCO Sanjeevani:} A dedicated family wellness and child care app designed specifically for employees and associated families, providing comprehensive health resources and support.
    
    \item \textbf{NISARG:} (Nalco Initiatives for Social Awareness and Responsible Growth) A mobile app that highlights the company's Corporate Social Responsibility (CSR) efforts, showcasing social impact projects and annual reports.
    
    \item \textbf{Hamesha Nalconian:} A dedicated portal for retired NALCO employees, serving as a single point for information on post-retirement medical reimbursements, empaneled hospitals, and related corporate policies.
    
    \item \textbf{NAGINAA:} (Nalco Grahak Information and Networking App for All) A customer-focused application providing details about products, pricing, MOU schemes, and dispatch information.
    
    \item \textbf{NAMASYA:} (Nalco Micro And Small enterprise Yogayog Application) An application tailored for MSE vendors to access registration processes, technical specifications, and vendor development programs.
    
    \item \textbf{Nalco Suraksha:} A mobile app utilized for on-site safety inspection reporting, instantly recording fire hazards, first aid requirements, and unsafe conditions at plant locations.
    
    \item \textbf{NALCO Smart Township:} An application developed for NALCO township residents and employees to raise property maintenance calls and seamlessly manage the Gate Pass system.
\end{itemize}